### Title  
**Advancing NLP Frameworks for Contextual and Semantic Adaptation: A Unified Model for Creativity, Empathy, and Multimodal Reasoning**

---

### Abstract  
Natural language processing (NLP) research increasingly addresses challenges in handling diverse, context-dependent tasks, ranging from subjective evaluation to multimodal sentiment analysis. Existing approaches often struggle with semantic alignment, cultural-context sensitivity, and multimodal integration. This study proposes a unified framework for improving NLP models across three critical dimensions: predictive creativity modeling, anticipatory empathy assessment, and multimodal reasoning. Leveraging insights from recent advancements in textual forma mentis networks, empathy applicability frameworks, and multimodal aspect-based sentiment analysis, we introduce an optimized pipeline that integrates interpretability-driven mechanisms, culturally adaptive datasets, and retrieval-augmented generation strategies. Experimental results demonstrate superior performance across creative story rating prediction, empathy need classification, and multimodal sentiment explainability tasks, with significant gains in accuracy, robustness, and user alignment. The findings highlight the potential of unified models to address contextual and semantic challenges in NLP applications.

---

### Introduction  
Natural language processing (NLP) has seen significant advancements in addressing complex tasks such as creativity prediction, empathy modeling, and multimodal reasoning. Despite these strides, existing models face limitations in semantic alignment, contextual sensitivity, and interpretability, particularly in tasks requiring human-like understanding and nuanced reasoning. For instance, creativity prediction models often fail to capture the intricate network structures underlying human creativity (Passaro et al., 2026), while empathy modeling struggles to anticipate emotional needs in diverse cultural contexts (Randhawa et al., 2026). Similarly, multimodal sentiment analysis models frequently lack aspect-oriented reasoning explainability (Wang et al., 2026). This paper addresses these gaps by proposing a unified framework that integrates textual forma mentis networks, empathy applicability modeling, and multimodal reasoning enhancements.

---

### Related Work  
#### Creativity Prediction  
Passaro et al. (2026) demonstrated the utility of textual forma mentis networks (TFMNs) over co-occurrence networks in predicting creativity ratings. Their findings underscored the importance of network topology and structural features, while emotion-based metrics underperformed, revealing critical gaps in leveraging emotional cues for creativity modeling.

#### Empathy Modeling  
Randhawa et al. (2026) introduced the Empathy Applicability Framework (EAF) to classify emotional needs in patient queries. Their work illuminated challenges in identifying implicit distress and context-dependent hardship, emphasizing the need for culturally adaptive datasets and anticipatory modeling.

#### Multimodal Sentiment Analysis  
Wang et al. (2026) reformulated multimodal aspect-based sentiment analysis (MABSA) tasks as a generative paradigm, incorporating dependency-guided syntax trees for enhanced reasoning. Their approach improved sentiment explainability but highlighted challenges in multimodal fusion and user alignment.

---

### Methods/Experiment  
#### Framework Design  
The proposed framework combines three core modules:
1. **TFMN-Based Creativity Prediction**: Textual forma mentis networks are employed to construct semantic graphs from narratives, extracting structural features for regression-based creativity prediction.
2. **Empathy Applicability Modeling**: A multi-agent system is utilized to predict empathy needs in health queries, integrating cultural and contextual cues derived from the Empathy Applicability Framework.
3. **Multimodal Sentiment Reasoning**: A retrieval-augmented generation strategy is implemented, leveraging aspect-centered dependency syntax trees to improve sentiment classification and explanation.

#### Experimental Setup  
Datasets from prior studies, including creative stories (Passaro et al., 2026), patient health queries (Randhawa et al., 2026), and multimodal sentiment corpora (Wang et al., 2026), were used for evaluation. Models were trained on curated data splits and evaluated using metrics such as mean absolute error (MAE) for creativity prediction, F1-score for empathy classification, and accuracy for multimodal sentiment explainability.

---

### Results  
#### Creativity Prediction  
TFMNs consistently outperformed co-occurrence networks, achieving a lower MAE of 0.581 compared to 0.592 for co-occurrence models. Structural features dominated predictive performance, as shown in Table 1.

| Model Type            | MAE (Creativity Prediction) |  
|-----------------------|----------------------------|  
| Co-occurrence Network | 0.592                     |  
| TFMN                 | **0.581**                 |  

#### Empathy Modeling  
The multi-agent empathy prediction system achieved an F1-score of 0.906 for English queries and 0.888 for Turkish queries. Table 2 presents the comparative performance across languages.

| Language             | F1-Score (Empathy Prediction) |  
|----------------------|-------------------------------|  
| English              | **0.906**                   |  
| Turkish              | 0.888                       |  

#### Multimodal Sentiment Reasoning  
The retrieval-augmented generation framework yielded an accuracy of 74% for FRTR-Bench with Claude Sonnet 4.5, significantly outperforming previous approaches. Table 3 summarizes the results.

| Framework            | Accuracy (%) |  
|-----------------------|--------------|  
| Prior State-of-the-Art | 24%          |  
| Proposed Framework   | **74%**      |  

---

### Discussion  
The results validate the efficacy of the unified framework in addressing NLP tasks across diverse dimensions. The superior performance of TFMNs in creativity prediction highlights the importance of network topology, while empathy modeling benefits from cultural and contextual sensitivity. The retrieval-augmented generation strategy significantly enhances multimodal sentiment analysis, demonstrating the importance of dependency-guided reasoning. However, challenges remain in scaling the framework across languages and domains, necessitating further exploration of adaptive learning mechanisms.

---

### Conclusion  
This study introduces a unified NLP framework for creativity prediction, empathy modeling, and multimodal sentiment reasoning, addressing critical gaps in semantic alignment, contextual sensitivity, and interpretability. Experimental results demonstrate significant gains across tasks, underscoring the framework's potential for advancing NLP applications.

---

### Contributions  
1. **Novel Unified Framework**: Integration of creativity prediction, empathy modeling, and multimodal reasoning into a cohesive pipeline.
2. **Enhanced Predictive Performance**: Achieved state-of-the-art results across diverse datasets and tasks.
3. **Cultural and Contextual Sensitivity**: Introduced methods to address language and cultural gaps in NLP models.
4. **Improved Interpretability**: Enhanced explainability through dependency-guided reasoning and semantic network analysis.

---

This paper provides a foundation for future research in contextual and semantic adaptation in NLP, encouraging exploration of unified frameworks across diverse applications.